\documentclass[standalone]{standalone}
\usepackage{tikz}
\usetikzlibrary{arrows.meta, calc, angles, quotes, positioning}

% 为了支持中文字符，需要使用能够处理中文的 LaTeX 引擎，
% 如 XeLaTeX，并在导言区加上以下宏包：
% \usepackage[UTF8]{ctex}

\begin{document}
	
	\begin{tikzpicture}[>=stealth, scale=1.3]
		
		% =========================================================================
		% 左图：惯性坐标系 (Inertial Frame / NED)
		% =========================================================================
		\begin{scope}[shift={(0,0)}]
			% 定义原点
			\coordinate (Oi) at (0, 0);
			
			% 绘制 x0 (North), y0 (East) 轴 (使用3D投影感)
			\draw[->, thick, gray] (Oi) -- (1.5, 1) coordinate (X0) node[above, black] {$x_0$};
			\draw[->, thick, gray] (Oi) -- (2.5, 0) coordinate (Y0) node[right, black] {$y_0$};
			
			% 绘制 z0 (Down) 轴 (垂直向下)
			\draw[->, thick, gray] (Oi) -- (0, -2) coordinate (Z0) node[below, black] {$z_0$};
			
			% 标注原点
			\node[left=2pt of Oi] {$o_0$};
			
			% 标题
			\node[below=15pt of Z0, font=\bfseries] {惯性坐标系};
		\end{scope}
		
		% =========================================================================
		% 中图：附体坐标系 (Body-fixed Frame) 6-DOF (右手定则修改版)
		% =========================================================================
		\begin{scope}[shift={(3.5,0)}]
			% 船体几何定义 (3D斜视图)
			\coordinate (SternTip) at (-1.5, -0.4);
			\coordinate (SternL)   at (-1.2, -0.6);
			\coordinate (SternR)   at (-1.2, -0.2);
			\coordinate (BowTop)    at (2.5,  0.5);
			\coordinate (BowR)     at (2.0,  0.8);
			\coordinate (BowL)     at (1.0,  0.0);
			\coordinate (BowTip)   at (2.8,  0.6);
			
			% 绘制船体轮廓并填充 (浅灰色)
			\draw[thick, fill=gray!15, join=round] 
			(SternTip) to[out=40, in=-140] (SternR) 
			-- (BowR) -- (BowTip) -- (BowL) -- (SternL)
			to[out=-40, in=140] (SternTip) -- cycle;
			
			% 定义体轴原点 o (在船体几何中心附近)
			\coordinate (Ob) at (1.0, 0.2);
			\filldraw[black] (Ob) circle (1.5pt) node[above left=1pt] {$o$};
			
			% --- 定义并标注 $x$ 轴 (向前 / Surge) ---
			\draw[->, ultra thick, blue!80!black] (Ob) -- (3.5, 0.7) coordinate (Xb) node[right] {$x$};
			% Surge 平移标注
			\draw[<->, thick, blue!80!black] ($(Ob)!0.4!(Xb)$) -- ++(0.5, 0.1) node[midway, sloped, above, font=\footnotesize] {纵荡};
			
			% --- 定义并标注 $y$ 轴 (向右 / 右手定则 / Sway) ---
			% 这部分需要绘制在船体下方，以体现3D空间感
			\draw[->, ultra thick, blue!80!black] (Ob) -- (2.0, -1.2) coordinate (Yb) node[below] {$y$};
			% Sway 平移标注
			\draw[<->, thick, blue!80!black] ($(Ob)!0.5!(Yb)$) -- ++(0.3, -0.2) node[midway, sloped, above, font=\footnotesize] {横荡};
			
			% --- 定义并标注 $z$ 轴 (向下 / Heave) ---
			\draw[->, ultra thick, blue!80!black] (Ob) -- (1.0, -1.8) coordinate (Zb) node[below] {$z$};
			% Heave 平移标注
			\draw[<->, thick, blue!80!black] ($(Ob)!0.5!(Zb)$) -- ++(0, -0.5) node[midway, right, font=\footnotesize] {垂荡};
			
			% --- 绘制 3D 旋转箭头 (Yaw, Pitch, Roll) ---
			
			% 1. 艏摇 (Yaw r) - 绕 $z$ 轴
			\draw[<-, thick, orange!80!black] (Ob) ++(-0.2, -0.6) arc (160:-160:0.3cm and 0.8cm);
			\node[orange!80!black, font=\small, text width=1cm, align=left] at (1.3, -1.2) {艏摇};
			
			% 2. 纵摇 (Pitch q) - 绕 $y$ 轴
			% 使用右手定则从 $z$ 转到 $x$，定义弧线
			\draw[->, thick, red!80!black] (2.2, -0.5) arc (-90:0:0.4cm and 0.2cm);
			\node[red!80!black, font=\small] at (2.3, -0.3) {纵摇};
			
			% 3. 横摇 (Roll p) - 绕 $x$ 轴 (在船首展示)
			\draw[->, thick, green!70!black] (BowTip) ++(-0.2, -0.2) arc (-120:120:0.25cm and 0.5cm);
			\node[green!70!black, font=\small] at (3.2, 0.3) {横摇};
			
			% 标题
			\node[below=15pt of Zb, font=\bfseries] {附体坐标系};
		\end{scope}
		
		% =========================================================================
		% 右图：平面投影与航向角 ψ
		% =========================================================================
		\begin{scope}[shift={(8,0)}]
			% 定义大地系平面坐标
			\draw[->, thick, gray] (0,-1) -- (0,3) coordinate (Y0p) node[above, black] {$x_0$};
			\draw[->, thick, gray] (-0.5,0) -- (4,0) coordinate (X0p) node[right, black] {$y_0$};
			\node[below left=2pt] at (0,0) {$o_0$};
			
			% 船体简化平面几何 (旋转 ψ 角度)
			\def\psiDeg{-35} % 顺时针旋转，所以是负值
			\coordinate (Pp) at (2.2, 1.5);
			\filldraw[black] (Pp) circle (1.5pt);
			
			\begin{scope}[shift={(Pp)}, rotate=\psiDeg]
				\draw[thick, fill=gray!10] (-0.7, -0.3) -- (0.5, -0.3) -- (0.8, 0) -- (0.5, 0.3) -- (-0.7, 0.3) -- cycle;
				
				% 绘制体轴 x, y 投影
				\draw[->, thick, blue!80!black] (0,0) -- (1.5, 0) coordinate (Xb_p) node[right] {$x$};
				\draw[->, thick, blue!80!black] (0,0) -- (0, -1.0) coordinate (Yb_p) node[below] {$y$};
			\end{scope}
			
			% 标注参考北向
			\draw[dashed, gray] (Pp) -- ++(0, 1.5) coordinate (N_ref);
			
			% 标注航向角 \psi (右手定则，顺时针为正)
			\pic [draw, ->, "$\psi$", angle eccentricity=1.3, angle radius=1.0cm, thick, green!60!black] 
			{angle = N_ref--Pp--Xb_p};
			
			% 标题
			%\node[below=15pt of (2,-1), font=\bfseries] {平面投影};
		\end{scope}
		
		% =========================================================================
		% 图表标题
		% =========================================================================
%		\node[below=45pt of {(Ob)}, font=\Large, text width=10cm, align=center] {图 3-1 \ 坐标系统示意图 (右手定则修正版)};
%		
	\end{tikzpicture}
	
\end{document}